\chapter{引用与注释}\label{chap:reference}

本章介绍参考文献引用、脚注和交叉引用的使用方法。

\section{参考文献引用}

在 \LaTeX{} 中用 \verb|\cite| 命令引用参考文献。参考文献格式应符合国标 GB/T 7714，本模板提供两种实现方式：
\begin{enumerate}
    \item 使用 \BibTeX{} 配合 \pkg{gbt7714} 宏包
    \item 使用 \BibLaTeX{} 配合 \pkg{biblatex-gb7714-2015} 样式包（其 \texttt{gb7714-2025} 样式实现 GB/T 7714-2025）
\end{enumerate}
其中，\BibLaTeX{} 方案功能更为丰富（支持 \cs{fullcite}、\cs{footfullcite} 等命令），且以 \texttt{biber} 为后端，对 Unicode 和中文文献支持更好，是本模板的\textbf{默认推荐}方案。不过，\texttt{biber} 的编译速度较 \texttt{bibtex} 慢，在文献条目较多时差异尤为明显。\BibTeX{} 方案编译更快，且部分期刊投稿仍要求使用，如对编译速度敏感或有投稿需求，可设置 \texttt{biblatex=false} 切换至 \BibTeX{} 方案。

无论选择哪种方式，均使用 \cs{tjbibresource} 命令在导言区指定参考文献数据库文件，使用 \cs{makereferences} 命令在正文中输出参考文献列表：

\begin{Verbatim}
\tjbibresource{bib/note.bib}        % 导言区指定 .bib 文件
...
\makereferences                     % 正文中输出参考文献
\end{Verbatim}

如需指定多个 \texttt{.bib} 文件，可用逗号分隔：\verb|\tjbibresource{file1.bib,file2.bib}|。

在正文中，使用 \cs{cite} 命令引用参考文献；\verb|key| 为参考文献数据库中条目的唯一标识符，多个键用逗号分隔：

\begin{Verbatim}
这是一个引用示例 \cite{book1}。
多文献引用 \cite{book1,online1}。
\end{Verbatim}

\cs{cite} 命令产生上标数字引用（如“……方法\textsuperscript{[1,2]}”），适合不希望打断行文的场合。

多个键可以一次引用，如\cite{book1,online1,article1}。以下按文献类型给出 \verb|\cite| 的引用示例：
\begin{itemize}
  \item 普通图书\cite{book1,book2}
  \item 论文集、会议录\cite{conf1,conf2}
  \item 科技报告\cite{techreport1,techreport2}
  \item 学位论文\cite{thesis1,thesis2,thesis3}
  \item 专利文献\cite{patent1,patent2}
  \item 专著中析出的文献\cite{inbook1,inbook2}
  \item 期刊中析出的文献\cite{qin2021,article1,guo_deepseek-r1_2025}
  \item 报纸中析出的文献\cite{newspaper1,newspaper2}
  \item 电子文献\cite{online1,online2,online3}
  \item 预印本\cite{preprint1,preprint2}
\end{itemize}

可使用 \verb|\nocite{key1,key2,key3...}| 将参考文献条目加入文献表中，但不在正文中引用。使用 \verb|\nocite{*}| 可将参考文献数据库中的所有条目加入文献表中。

在编写 \texttt{.bib} 文件时，需注意以下常见问题：
\begin{itemize}
  \item 每个条目的键（如 \texttt{book1}）在整个 \texttt{.bib} 文件中必须唯一；
  \item 常见的条目类型包括 \texttt{@book}（图书）、\texttt{@article}（期刊论文）、\texttt{@inproceedings}（会议论文）、\texttt{@phdthesis}（博士论文）、\texttt{@mastersthesis}（硕士论文）、\texttt{@online}（网络资源）、\texttt{@preprint}（预印本，GB/T 7714-2025 新增）等；
  \item 中文作者姓名之间使用 \texttt{and} 连接，而非顿号或逗号；
  \item 撰写规范要求参考文献中的标点统一采用半角，样式只能控制它生成的标点；题名等字段里自带的标点请直接写成半角，引号写作 \verb|``…''|，括号写作 \verb|(…)|；
  \item 版本项只写数字，如 \verb|edition = {3}|，由样式排成“3rd ed.”；
  \item 英文文献的题名与书名按句首大写输出，期刊名保持原样；需要保持原有大小写的专有名词与缩写须用花括号括起，如 \verb|{IEEE}|、\verb|{PyTorch}|；网址写在 \texttt{url} 字段，不要写进题名；
  \item 汉语拼音姓名的名按全称输出（如“He Kaiming”），西文姓名的名缩写为首字母（如“Knuth D E”）；拼音名须写成一个词，如 \verb|He, Kaiming|，两个音节之间可用撇号隔开（\verb|Xi'an|），写成 \verb|Kai-Ming| 或 \verb|Kai Ming| 会被当作西文名缩写；拼音姓不要加花括号；
  \item 建议使用 Zotero、Mendeley 等软件管理参考文献，可自动生成规范的 \texttt{.bib} 文件。
\end{itemize}

% \fullcite 仅在 biblatex=true 时定义；未定义则跳过本节
\ifx\fullcite\undefined\else
\subsection{\BibLaTeX{}特有的引用命令}

使用 \pkg{biblatex} 时，还可以使用以下引用命令。

\subsubsection{完整引用}

需要在正文中（参考文献章节以外）插入某条文献的完整引用时，可以使用 \verb|\fullcite{key1}| 命令，示例如下：

\begin{itemize}
    \item \fullcite{qin2021}
\end{itemize}

\subsubsection{脚注引用}

以脚注形式引用文献时，可以使用 \verb|\footfullcite{key1}| 命令\footfullcite{qin2021}。
\fi


\section{脚注}

脚注是排在页面底部的注释或补充说明。在 \LaTeX{} 中，在需要加注的词语或句子后面使用 \verb|\footnote| 命令即可，如下所示：

\begin{Verbatim}
脚注是排在页面底部的注释或补充说明\footnote{通常，我们在脚注里也写完整的句子。在文本中使用脚注时，应该遵循学术规范，尽可能引用可信的来源，并注明出处。}。
\end{Verbatim}

其中，花括号中的文本就是脚注的内容。编译文档后，脚注会出现在页面底部，并自动标上数字：脚注是排在页面底部的注释或补充说明\footnote{通常，我们在脚注里也写完整的句子。在文本中使用脚注时，应该遵循学术规范，尽可能引用可信的来源，并注明出处。}。

按《同济大学本科生毕业设计（论文）撰写规范》，脚注编号以 \TJCircled{1}、\TJCircled{2}、\TJCircled{3} 等带圈形式呈现，由模板自动生成，无需手动设置。

\section{交叉引用}

交叉引用指在文中引用章节、图表、公式或页码等对象。在 \LaTeX{} 中，用 \verb|\label| 命令为对象添加标签，再用 \verb|\ref| 或 \verb|\pageref| 命令引用。

为便于管理，建议为标签采用统一的命名前缀：
\begin{center}
\begin{tabular*}{\linewidth}{@{\extracolsep{\fill}}llll}
  \toprule
  章：\verb|chap:a| & 节：\verb|sec:b| & 图：\verb|fig:c| & 表：\verb|tab:d| \\
  公式：\verb|eq:e| & 算法：\verb|algo:f| & 代码：\verb|lst:g| & 定理：\verb|thm:h| \\
  \bottomrule
\end{tabular*}
\end{center}

例如，在文档中的某个章节中添加一个标签：

\begin{Verbatim}
\chapter{引言}
\label{chap:intro}
这是一段引言。
\end{Verbatim}

然后，在文档的其他位置使用 \verb|\ref| 命令来引用这个标签：

\begin{Verbatim}
请参见第 \ref{chap:intro} 章。
\end{Verbatim}

这样，就可以在文档中产生“请参见第 X 章”的效果，其中 X 为标签所在章的编号。

类似地，使用 \verb|\pageref| 命令可以引用页面编号。例如：

\begin{Verbatim}
请参见第 \pageref{chap:intro} 页。
\end{Verbatim}

这样，就可以在文档中产生“请参见第 X 页”的效果，其中 X 为标签所在页的页码。

为免于手动维护“第 X 章 / 图 X / 表 X”等引用词汇，本模板加载了 \pkg{cleveref} 宏包，其 \cs{cref} 命令会根据标签所指对象的类型自动生成引用词汇：

\begin{Verbatim}
请参见\cref{chap:intro}。   % 产出：请参见第 1 章。
\end{Verbatim}

本模板已用 \cs{crefformat} 为章、节、图、表、公式、算法、代码和各定理类环境设置了引用格式，\cs{crefformat} 的设置优先于 \cs{crefname}，所以自定义引用词汇时应在导言区重新声明 \cs{crefformat}，其中 \verb|#1| 为编号，\verb|#2| 与 \verb|#3| 标出超链接的起止；一次引用多个同类对象时使用的格式由 \cs{crefmultiformat} 设置。以编号结尾的格式在 \verb|#3| 之后加 \cs{tjrefnum}，多项格式在“和”之前加 \cs{tjrefsep}，XeLaTeX 才会在编号与后面的汉字之间留出间距：

\begin{Verbatim}
\crefformat{对象类型}{#2引用词汇~#1#3\tjrefnum}
\end{Verbatim}

定理类环境的对象类型就是环境名，如定理为 \texttt{thm}、推论为 \texttt{cor}。例如，可以使用以下代码将定理的引用词汇修改为“命题”：
\begin{Verbatim}
\crefformat{thm}{#2命题~#1#3\tjrefnum}
\crefmultiformat{thm}{#2命题~#1#3\tjrefnum}{\tjrefsep 和#2命题~#1#3\tjrefnum}%
  {、#2命题~#1#3\tjrefnum}{\tjrefsep 和#2命题~#1#3\tjrefnum}
\end{Verbatim}

需要同时引用几个并列的对象时，可以直接使用 \verb|\cref| 命令，例如：
\begin{Verbatim}
请参见\cref{chap:introduction,chap:conclusion}。
请参见\cref{chap:introduction,chap:math,chap:reference,chap:float}。
请参见\cref{chap:introduction,chap:reference,chap:conclusion}。
请参见\cref{chap:introduction,chap:reference,chap:conclusion,%
  algo:algorithm,fig:parallel1,tab:firstone}。
\end{Verbatim}

引用多个对象时，\pkg{cleveref} 会自动合并连续编号为范围，并按类型分组显示：
\begin{itemize}
  \item 请参见\cref{chap:introduction,chap:conclusion}。
  \item 请参见\cref{chap:introduction,chap:math,chap:reference,chap:float}。
  \item 请参见\cref{chap:introduction,chap:reference,chap:conclusion}。
  \item 请参见\cref{chap:introduction,chap:reference,chap:conclusion,%
        algo:algorithm,fig:parallel1,tab:firstone}。
\end{itemize}
